% GENERATED_BY: oc_core_monograph_translation_draft_pipeline_v1.py
% AI_ASSISTED_TRANSLATION_DRAFT: TRUE
% THEOREM_REVIEW_STATUS: PENDING
% THEOREM_REVIEW_MODE: FORMAL_AI_GATE__CODEX_CLI_CHATGPT
% THEOREM_REVIEW_REVIEWER_ID:
% THEOREM_REVIEW_AUTHORITY_CLASS:
% THEOREM_REVIEWED_AT:
% THEOREM_REVIEW_DOSSIER_REF:
% SOURCE_LANGUAGE: EN
% TARGET_LANGUAGE: RU
% SOURCE_EN_REF: content/k_levels/k0.tex
% ================================================================
% ==== FILE: content/k_levels/k0.tex
% ================================================================

% ==============================
%  Ontology of Continua — Core
%  K-level Module: K0
%  File: content/k_levels/k0.tex
%  Status: FULLY DEFINED
% ==============================

\subsubsection{\texorpdfstring{$K_0$}{K_0} Обзор}
\label{sec:k0-overview}

Уровень $K_0$ представляет собой предпредструктурный фундамент всех континуумов.
Он не является физической или геометрической областью, а логической
подосновой, определяющей
\emph{условия возможности} для любого континуума $K_x$.
Он обеспечивает метадавление, ограничения допустимости, логическое
превосходство (аксиома~0.3) и универсальные правила, определяющие,
какие структуры могут существовать.

$K_0$ сам по себе не обладает обычными осями или геометрическими
размерностями.
Вместо этого он задаёт:
\begin{itemize}
    \item мета-аксиомы $\mathcal{L}$, определяющие структурную согласованность,
    \item пространство возможностей $\mathcal{P}$ для всех континуумов,
    \item метадавление, порождающее возникновение различий,
    \item универсальные пороги, задающие допустимость и коллапс,
    \item правила рекурсии $K_x \mapsto M_x \mapsto K_{x+1}$.
\end{itemize}

Его главная функция —
\textbf{логическое управление} всеми более глубокими уровнями.

% ================================================================
\subsubsection{\texorpdfstring{$\Omega(K_0)$}{\Omega(K_0)} Пространство состояний}

$\Omega(K_0)$ состоит из всех допустимых конфигураций пары
\[
(\mathcal{P}, \mathcal{L}),
\]
где $\mathcal{P}$ — это метадомен возможных структур, а
$\mathcal{L}$ — система мета-аксиом, обеспечивающая согласованность.

Поскольку $K_0$ неметричен и негеометричен, $\Omega(K_0)$ не является
метрическим пространством, а является абстрактным множеством со следующими
свойствами:
\begin{enumerate}
    \item $\Omega(K_0)$ содержит все возможные спецификации структурных
          ограничений, позволяющих существование континуумов.
    \item У $\Omega(K_0)$ нет внутренней топологии, кроме той,
          которая индуцирована ограничениями в $\mathcal{L}$.
    \item Для каждого $\Omega(K_x)$, где $x \ge 1$, выполняется

          \[
              \Omega(K_x) \subseteq \Omega(K_0).
          \]
\end{enumerate}

Внутри $\Omega(K_0)$ нет непрерывной эволюции: он задаёт
\emph{условия}, а не динамические траектории.

% ================================================================
\subsubsection{\texorpdfstring{$\partial\Omega(K_0)$}{\partial\Omega(K_0)} Граница}

Граница $K_0$ включает все метаконфигурации, нарушающие
внутренние условия согласованности мета-аксиом $\mathcal{L}$.

Конфигурация $(\mathcal{P},\mathcal{L})$ принадлежит
$\partial\Omega(K_0)$, если:
\begin{itemize}
    \item аксиомы противоречат друг другу,
    \item допустимость схлопывается ($\mathcal{P}$ становится пустым),
    \item правила рекурсии становятся неопределёнными,
    \item структурные условия согласованности не поддерживаются.
\end{itemize}

Переход границы $\partial\Omega(K_0)$ соответствует логической
невозможности существования
любого континуума.

% ================================================================
\subsubsection{\texorpdfstring{$A(K_0)$}{A(K_0)} Оси}

$K_0$ не имеет геометрических осей. Вместо этого в нём присутствуют:
\begin{itemize}
    \item абстрактные оси
\emph{возможности},
    \item оси допустимости, управляющие тем, какие структуры
          могут быть получены,
    \item мета-оси, задающие логические зависимости между уровнями.
\end{itemize}

В формальном виде:
\[
A(K_0) = \{ \text{размерности допустимых структурных вариаций} \}.
\]

Эти оси — исходный источник всех последующих осей:
\[
A(K_x) \subseteq A(K_0).
\]

% ================================================================
\subsubsection{\texorpdfstring{$P(K_0)$}{P(K_0)} Потенциалы}

Потенциалы в $K_0$ описывают мета-условия, побуждающие континуумы к
формированию структуры:
\[
P(K_0) = \{ \text{напряжения, ограничения, логические градиенты} \}.
\]

Они включают:
\begin{itemize}
    \item метадавление (аксиома 0.z),
    \item давление к структуризации,
    \item давление к совместимым различиям,
    \item давление к непустым допустимым пространствам.
\end{itemize}

Эти потенциалы не являются числовыми энергиями, а логическими
интенсивностями.

% ================================================================
\subsubsection{\texorpdfstring{$\Theta(K_0)$}{\Theta(K_0)} Пороги}

Пороги в $K_0$ соответствуют логическим пороговым значениям:
\[
\Theta(K_0) = \{ \Theta_{\mathrm{coh}}, \Theta_{\mathrm{adm}},
                 \Theta_{\mathrm{meta}} \}.
\]

Примеры:
\begin{itemize}
    \item $\Theta_{\mathrm{coh}}$: минимальная согласованность мета-аксиом,
    \item $\Theta_{\mathrm{adm}}$: минимальное непустое пространство
          возможности,
    \item $\Theta_{\mathrm{meta}}$: минимальная поддержка рекурсии.
\end{itemize}

Нарушение любого порога делает невозможным существование континуумов на
любом $K_x$.

% ================================================================
\subsubsection{\texorpdfstring{$J(K_0)$}{J(K_0)} Потоки}

В $K_0$ отсутствуют физические потоки.
$J(K_0)$ включает:
\begin{itemize}
    \item потоки допустимости,
    \item потоки логической зависимости,
    \item распространение ограничений,
    \item метапотоки, описывающие, какие структуры могут порождать
          другие.
\end{itemize}

Эти потоки формируют пространство возможных континуумов, но не
соответствуют эволюции во времени.

% ================================================================
\subsubsection{\texorpdfstring{$C(K_0)$}{C(K_0)} Циклы}

Циклы в $K_0$ соответствуют:
\begin{itemize}
    \item циклам фиксированных точек согласованности,
    \item мета-устойчивым петлям допустимости,
    \item самореферентным циклам замыкания в $\mathcal{L}$.
\end{itemize}

Цикл существует, если:
\[
\Pi(C) > \Theta_{\mathrm{coh}},
\]
где $\Pi(C)$ — это «мощность» мета-цикла (устойчивость рекурсивного
замыкания).

В $K_0$ нет временных циклов; это лишь логические циклы.

% ================================================================
\subsubsection{\texorpdfstring{$\tau(K_0)$}{\tau(K_0)} Время}

Время не существует в $K_0$.

Теорема «Возникновение времени из C-циклов» (Physics Run):

Время возникает, когда $\Omega(K_2)$ содержит неразрушаемый цикл
$C$ такой, что:
\[
\Pi(C) > \Theta_{\mathrm{time}}.
\]

Следовательно:
\[
\tau(K_0) \;\text{не определено}.
\]

% ================================================================
\subsubsection{\texorpdfstring{$k(K_0)$}{k(K_0)} Континуумность}

Континуумность измеряет допустимость всех континуумов:
\[
k(K_0) =
H(\Omega(K_0)) \cdot
k_A(K_0) \cdot
k_{\mathrm{coh}}(K_0),
\]
где:
\begin{itemize}
    \item $H(\Omega(K_0)) = 1$, если $\Omega(K_0)\neq\emptyset$, иначе $0$,
    \item $k_A(K_0)$ измеряет богатство допустимых осей,
    \item $k_{\mathrm{coh}}$ измеряет согласованность мета-аксиом.
\end{itemize}

Если $k(K_0)=0$, ни один континуум на более высоком уровне не может
существовать.

% ================================================================
\subsubsection{\texorpdfstring{$T(K_0)$}{T(K_0)} Структурное напряжение}

Структурное напряжение определяется как:
\[
T(K_0) = T(\mathcal{P}, \mathcal{L}),
\]
представляя степень несовместимости внутри мета-аксиом или между
мета-аксиомами и допустимыми структурами.

Высокое напряжение может:
\begin{itemize}
    \item уменьшать допустимые множества,
    \item порождать несовместимые ветви,
    \item приводить к коллапсу $\Omega(K_0)$.
\end{itemize}

% ================================================================
\subsubsection{\texorpdfstring{$E(K_0)$}{E(K_0)} Энергия}

В $K_0$ отсутствует физическая энергия.
Вместо этого:
\[
E(K_0) = E_{\mathrm{meta}},
\]
формальная мера внутренней поддержки для согласованных структур:
\[
E_{\mathrm{meta}} \propto \text{согласованность }\mathcal{L}.
\]

% ================================================================
\subsubsection{\texorpdfstring{$K_0$}{K_0} Операторы (\texorpdfstring{$\Psi$}{\Psi}, \texorpdfstring{$\Phi$}{\Phi}, \texorpdfstring{$\Lambda$}{\Lambda}, \texorpdfstring{$U$}{U}, \texorpdfstring{$\Chi$}{\Chi})}

Операторы в $K_0$ управляют допустимостью и рекурсией:

\begin{itemize}
    \item $\Psi$: оператор порождения уровней

          \[
          \Psi: K_x \mapsto M_{x+1}
          \]
    \item $\Phi$: оператор эволюции мета-пространства

          \[
          \Phi: M_x \mapsto M_x(t+dt)
          \]
    \item $\Lambda$: оператор ограничения, вводящий $\mathcal{L}$
    \item $U$: оператор структурной унификации
    \item $\Chi$: оператор логического ветвления (аксиома~21)
\end{itemize}

Эти операторы определяют каркас вертикальной рекурсии.

% ================================================================
\subsubsection{Процессы на \texorpdfstring{$K_0$}{K_0}}

Ключевые процессы:
\begin{itemize}
    \item логическое унифицирование,
    \item расширение/сжатие допустимости,
    \item мета-стабилизация рекурсивных определений,
    \item ветвление пространств возможностей,
    \item коллапс несовместимых мета-структур.
\end{itemize}

% ================================================================
\subsubsection{\texorpdfstring{$K_0$}{K_0} Прогнозы}

Теория предсказывает:
\begin{itemize}
    \item отсутствие наблюдаемой динамики в $K_0$,
    \item универсальность порогов для всех уровней,
    \item неизбежность размерностного роста под действием метадавления,
    \item существование метаонтологических структур ветвления.
\end{itemize}

% ================================================================
\subsubsection{\texorpdfstring{$K_0$}{K_0} Эксперименты}

$K_0$ невозможно экспериментально исследовать напрямую.
Косвенные проверки:
\begin{itemize}
    \item универсальность поведения порогов в физических,
          химических, биологических, когнитивных и социальных системах,
    \item распространённость структурных рекурсивных паттернов,
    \item фальсифицируемые тесты возникающего времени (Physics Run),
    \item устойчивость структур ветвления (K$_{11}$–K$_{12}$).
\end{itemize}

% ================================================================
\subsubsection{Смерть \texorpdfstring{$K_0$}{K_0}}

Коллапс наступает, если:
\[
\Omega(K_0) = \varnothing.
\]

Это соответствует:
\begin{itemize}
    \item противоречивым мета-аксиомам,
    \item логически несогласованному пространству возможностей,
    \item нерушимой рекурсии $\Psi$ или $\Phi$,
    \item нарушению $\Theta_{\mathrm{coh}}$ или $\Theta_{\mathrm{adm}}$.
\end{itemize}

Если $K_0$ коллапсирует, все континуумы коллапсируют одновременно.

% ================================================================
\subsubsection{\texorpdfstring{$K_0$}{K_0} Фальсифицируемость}

Следующие предсказания фальсифицируемы:
\begin{enumerate}
    \item Универсальное присутствие порогов $\Theta$ на всех уровнях.
    \item Необратимость размерностного роста (Теорема~1).
    \item Возникновение времени только на $K_2$ через циклы.
    \item Существование универсальной рекурсии $K_x\to M_x\to K_{x+1}$.
\end{enumerate}

Эмпирическое отрицание любого из этих утверждений фальсифицирует $K_0$.

% ================================================================
\subsubsection{\texorpdfstring{$K_0$}{K_0} Ветвление / Онтологическая позиция}

На основе аксиомы~21 (онтологическое ветвление) $K_0$ является корневым
узлом дерева онтологического ветвления:
\[
K_0 \rightarrow K_1 \rightarrow K_2 \rightarrow \dots \rightarrow K_{12}.
\]

Все ветви обязаны оставаться совместимыми с $\mathcal{L}$; несовместимые
ветви уничтожаются коллапсом допустимости.

% ================================================================
\subsubsection{\texorpdfstring{$K_0$}{K_0} Связь с M-пространствами}

$K_0$ задаёт допустимость и структуру всех метапространств:
\[
M_x \subseteq \Omega(K_0).
\]

И, напротив:
\[
K_x \subseteq M_x,
\qquad
M_x \subseteq \Omega(K_{x+1}),
\]
так что $K_0$ является единственным унифицирующим суперпространством для
всей иерархии.
